\section{Where Unit State Lives: The Three-Home State Model}
\label{sec:sec04}
\label{sec:states}

Unit state lives in three homes, one map each: immutable \emph{ProductTerms}$[u]$
keyed by unit; shared \emph{UnitStatus}$[u]$ keyed by unit; and per-position
\emph{PositionState}$[w,u]$ keyed by (wallet, unit). There is no wallet-keyed state
sector: every per-wallet economic fact is keyed by a mandate or strategy unit and so
lives in \emph{PositionState}$[w,u_{\mathrm{MA}}]$. Each economic fact has exactly one
home and one writer, so internal reconciliation failure cannot arise; conservation and
deterministic replay then hold by construction, through the sealed single-writer
discipline and the purity and totality of replay. This placement is the minimum basis of
the problem (\S\ref{sec:sec04-three}) and unique under any correctness gate
(\S\ref{sec:sec04-pareto}).

A \emph{unit} is anything held and conserved (\S\ref{sec:ledger}): an instrument, and
also a mandate or a strategy contract (\S\ref{sec:sec04-three}). The framework
primitives --- wallet, unit, holding $h(w,u)$, move, transaction, conservation
$\sum_{w}h(w,u)=0$ --- are those of \S\ref{sec:ledger}.

\subsection{The placement rule and the 2\texorpdfstring{$\times$}{x}2}
\label{sec:sec04-rule}

One rule governs every placement: each economic fact has exactly one home and one
writer. A fact violates it in exactly two ways. \emph{Keyed wrong}: one home collapses
two facts, or one fact duplicates across homes that drift. \emph{Authored wrong}: two
authorities keep one fact and their records drift apart. Two questions forestall the two
failures, and because there are only two failures there is no third question.

\emph{Does the fact depend on the holder, or only on the unit?} The scope is one unit's
state, so a unit and a holder of it are the only terms in view. Every economic fact ---
one entering conservation, valuation, or profit and loss --- is then a (holder, unit)
fact or a unit fact. What stays keyed by the wallet alone --- KYC, permissions, an audit
cursor --- enters none of these three reckonings; it is identity, not economic state, and
lives in a \emph{WalletRegistry} that carries no quantity and joins no conservation law.

\emph{Is the fact's record owned by an outside authority or by the ledger?} The criterion
is who owns the record's history, not who sources the number. An externally authored fact
is one an outside authority --- the exchange, the contract, the reference-data provider
--- owns and restates, and the ledger preserves version by version. A ledger-authored
fact is one the ledger's own events produce and overwrite. A benchmark level the ledger
records and overwrites is ledger-authored; the benchmark identity the provider restates is
the authority's versioned record --- the two sort apart though both come from that
provider.

The two questions give a $2\times2$; three of its four cells are occupied, and each
occupied cell is one home.

\begin{center}
\begin{tabular}{@{}r p{0.36\textwidth} p{0.36\textwidth}@{}}
\toprule
 & \textbf{ledger-authored} & \textbf{externally authored} \\
\midrule
\textbf{per unit} & \textbf{Status}: lifecycle stage, last settlement price, current
weights, benchmark level & \textbf{Terms}: multiplier, currency, expiry, strike, ISIN,
fee schedule, mandate text, benchmark identity \\
\addlinespace
\textbf{per (holder, unit)} & \textbf{Position}: held quantity, accumulated cost,
high-water mark, entry NAV & \textit{--- empty ---} \\
\bottomrule
\end{tabular}
\end{center}

Terms and status share the unit key but cannot be one value, since their authorities of
record differ (\S\ref{sec:sec04-three}); they ride together on the unit key as a
co-present pair --- two homes whose mutation disciplines differ, each kept in its own map
(\S\ref{sec:sec04-construction}).

\subsection{Why three homes, and why the fourth cell is empty}
\label{sec:sec04-three}

Each occupied cell exists for one reason; the empty cell is empty for one; and exactly
three maps are forced, no fewer and no more.

\emph{Position is keyed by (holder, unit), because two holders of one contract hold
different state.} A buyer and a seller hold the same future; the buyer's quantity is
$+1000$, the seller's $-1000$, and their accumulated costs differ. Substituting one
holder's position for the other's changes the value, so any unit-only keying identifies the
two and loses the difference. The holder is in the key.

\emph{Status is keyed by the unit alone, because one value is read identically by every
holder.} A future settles at one price on the day; an index stands at one level. Stored
per holder, the number is copied across thousands of rows free to drift --- a
reconciliation break by construction. Keyed by the unit, it is one value with one writer.

\emph{Terms are a home distinct from status, because the two have different authorities of
record.} Terms are externally authored, status ledger-authored. Co-mingling the
authority's record with the ledger's own is the single-source-of-truth violation the
system exists to prevent. Terms are append-only and versioned for reconstruction; the
status cache is overwritten on every settlement while the settling event is appended
immutably to the log (\S\ref{sec:sec04-unitstatus}). Merging the two would put two
mutation disciplines under one name.

\emph{The fourth cell is empty, because no authority issues a fact about one holder's
position.} A position is the ledger's own from origin: held quantity folded from the move
stream, high-water mark and entry NAV written by ledger events. A custodian statement or
prime-broker position report is a reconciliation input the ledger checks against, not a
record it adopts as an authority's own. So the (holder, unit) key carries only
ledger-authored facts, and the externally authored cell is empty by construction.

\paragraph{The $W$-sector collapses (C12).}
A wallet-keyed economic-state map --- call it the \emph{$W$-sector},
$\mathrm{Map}[\mathrm{WalletId},\mathrm{WalletState}]$ --- is the apparent demand of the
managed account, whose state is the client's, not the instrument's. It does not exist,
because the mandate is a unit. The manager issues a mandate unit $u_{\mathrm{MA}}$ to the
client (\S\ref{sec:managed}), conserving like any issued unit:
\[
h(w_{\mathrm{manager}},u_{\mathrm{MA}})=-1,\qquad
h(w_{\mathrm{client}},u_{\mathrm{MA}})=+1,\qquad
\sum_{w}h(w,u_{\mathrm{MA}})=0 .
\]
Every field that looked per-wallet --- high-water-mark value, accrued management and
performance fee, mandate breach flags, benchmark NAV at this wallet's inception --- is then
a fact about the position $(w_{\mathrm{client}},u_{\mathrm{MA}})$ and lives in
\emph{PositionState}. Two mandates on one wallet are two rows, each with its own
high-water mark and fees; a wallet-keyed scalar would have merged them. This is why
\hyperref[cond:C12]{C12} keys on the mandate unit, not the wallet. A relationship spanning
several instruments that does \emph{not} reduce to a single (holder, unit) row would be a
(holder, several-units) fact occupying a fourth home and a third map; the count excludes
it, and it lies outside this scope.

\subsection{The UnitStatus discipline, stated once}
\label{sec:sec04-unitstatus}

UnitStatus holds one value per unit, shared---read identically by every holder. Its
value changes over time, but UnitStatus is not a separate source of truth: the immutable
event log is. Every change is caused by a logged event; the stored value is overwritten in
place only as a read cache. Replaying the events up to any point rebuilds the exact value
that held then, so nothing is lost by overwriting and there is no other way to change the
value. A value exists from registration onward.

Equivalently, the stored value always equals what replaying the unit's events produces ---
the catamorphism property, meaning the fold of the event stream always rebuilds the exact
cached value --- and no part of the system writes it except through a logged event. In the
reference (\S\ref{sec:sec04-reference}) this holds by construction: every
\emph{UnitStatus} field is written only by \texttt{applyStatus}, folded over the unit's
status writes inside the sole transition \texttt{applyTx} (and \texttt{register}, which is
\texttt{applyTx} of a move-less introduction); the
\texttt{Ledger} constructor is unexported, so no out-of-band \texttt{setStatus} exists.
Externally sourced inputs --- a settlement price, the current benchmark level --- enter
only as a logged observation event, so the fold stays pure at the boundary. Every later
reference to cached unit state in this specification is read in this sense.

\subsection{The construction}
\label{sec:sec04-construction}

The homes are built from the smallest pieces, each forced by the one before, and the
disciplines C1--C12 (\S\ref{sec:sec04-conditions}) are introduced where the construction
forces them. Quantities are exact integer minor units, never floating point, so
determinism and conservation are arithmetic facts.

\paragraph{Scalars and keys.}
A balance is a quantity with a zero, associative addition, and a negation --- an abelian
group. The negation forces the two legs of a transfer to cancel; the zero makes the
sum over no wallets conserved with no special case, which discharges the vacuous base of
conservation (\hyperref[cond:C9]{C9}). A price is a number but not a quantity --- prices
are never summed --- so it is a separate type with neither identity nor inverse. A balance
is named by two keys, wallet and unit, the holder in the key.

\begin{lstlisting}[language=Haskell]
{-# LANGUAGE GADTs, DataKinds, KindSignatures, StandaloneDeriving #-}
import           Control.Monad      (foldM)
import           Data.List          (foldl')
import           Data.List.NonEmpty (NonEmpty (..))
import qualified Data.List.NonEmpty as NE
import           Data.Map.Strict    (Map)
import qualified Data.Map.Strict    as Map

newtype Qty = Qty Integer deriving (Eq, Ord)        -- exact minor units, never Float
instance Show Qty      where show (Qty n) = show n
instance Semigroup Qty where Qty a <> Qty b = Qty (a + b)
instance Monoid    Qty where mempty = Qty 0          -- the zero discharges the vacuous base (C9)

qneg :: Qty -> Qty                                   -- group inverse: a transfer's legs cancel
qneg (Qty n) = Qty (negate n)
qmax :: Qty -> Qty -> Qty
qmax (Qty a) (Qty b) = Qty (max a b)

newtype WalletId = WalletId String deriving (Eq, Ord, Show)
newtype UnitId   = UnitId   String deriving (Eq, Ord, Show)
\end{lstlisting}

\paragraph{ProductTerms --- the authority's record, grown only by appending (C6, C7).}
``One or more versions, append-only'' is a non-empty list, so a registered unit always has
a current version and ``registered but versionless'' is not representable. The type is
abstract: its constructor is unexported, so no importer can lay down a fresh one-version
value and discard history. The only growth operations are \texttt{register}, which writes
version one, and \texttt{appendVersion}, which appends; no setter exists, so the
append-only discipline holds by absence, not by convention (\hyperref[cond:C6]{C6}). The
first version is written at registration, so terms are total on the registered units
(\hyperref[cond:C7]{C7}).

\begin{lstlisting}[language=Haskell]
data TermsVersion = TermsVersion
  { tvLabel :: String, tvFields :: Map String String }  -- opaque payload: multiplier, ISIN, ...
  deriving (Eq, Show)

newtype ProductTerms = ProductTerms (NonEmpty TermsVersion) deriving (Show)
  -- exported ABSTRACT: no data constructor, no setter -> append-only holds by absence (C6)

currentTerms :: ProductTerms -> TermsVersion           -- total: NonEmpty always has a last
currentTerms (ProductTerms vs) = NE.last vs
allVersions  :: ProductTerms -> NonEmpty TermsVersion
allVersions  (ProductTerms vs) = vs
appendVersion :: TermsVersion -> ProductTerms -> ProductTerms   -- C6: append, never rewrite
appendVersion tv (ProductTerms vs) = ProductTerms (vs <> (tv :| []))

productTerms :: Ledger -> UnitId -> Maybe ProductTerms  -- Nothing = unregistered
productTerms l u = Map.lookup u (ledgerPT l)
\end{lstlisting}

\paragraph{UnitStatus --- the shared read cache (C5).}
Status records the unit's lifecycle stage together with its last settlement price and any
successor, under the discipline of \S\ref{sec:sec04-unitstatus}. The lifecycle is a closed
sum, \texttt{Registered}$\to$\texttt{Active}$\to$\texttt{Expired}$\to$\texttt{Closed};
\texttt{Registered} is the at-registration stage --- a unit recorded in and known to the
ledger before any settlement, which an OTC option, a bespoke mandate, or an internal
strategy reaches exactly as an exchange instrument does.
Every \emph{UnitStatus} field changes only through a \texttt{StatusWrite} --- the closed
writer set whose three constructors each name one field --- applied by the single total
writer \texttt{applyStatus}, the projection step of \S\ref{sec:sec04-unitstatus}; no other
code overwrites a status field, so each write is last-write-wins and idempotent (P6).
Product-declared defaults are written at Unit Store registration (\S\ref{sec:unit-store}),
so status is total on the registered units (\hyperref[cond:C5]{C5}).

\begin{lstlisting}[language=Haskell]
data Lifecycle = Registered | Active | Expired | Closed deriving (Eq, Show)

data UnitStatus = UnitStatus
  { usLifecycle    :: Lifecycle
  , usLastSettle   :: Maybe Qty       -- None at the registration default (C5)
  , usSupersededBy :: Maybe UnitId    -- set only by a Breaking amendment (C8)
  } deriving (Eq, Show)

defaultStatus :: UnitStatus           -- C5: product-declared default, written at registration
defaultStatus = UnitStatus Registered Nothing Nothing

data StatusWrite                      -- the closed status-writer set; each names ONE field
  = SetLifecycle    Lifecycle         -- a lifecycle event
  | SetLastSettle   Qty               -- a settle event (a logged observation)
  | SetSupersededBy UnitId            -- a Breaking amendment (C8)
  deriving (Eq, Show)

applyStatus :: StatusWrite -> UnitStatus -> UnitStatus   -- the ONLY writer of a status field;
applyStatus (SetLifecycle l)    s = s { usLifecycle    = l }      -- total, last-write-wins,
applyStatus (SetLastSettle q)   s = s { usLastSettle   = Just q } -- idempotent (P6)
applyStatus (SetSupersededBy u) s = s { usSupersededBy = Just u }

unitStatus :: Ledger -> UnitId -> Maybe UnitStatus       -- Nothing = unregistered
unitStatus l u = Map.lookup u (ledgerUS l)
\end{lstlisting}

\paragraph{PositionState --- the per-position home (C1).}
A position carries more than a balance: its primary field is the held quantity, which
conserves; beside it ride non-conserved per-(holder, unit) facts --- a high-water mark, an
entry NAV --- witnessing that the Position home carries facts no zero-sum law binds. A
first touch starts a position at $0_P$, the all-zero state. Two orthogonal disciplines
govern the home (\hyperref[cond:C1]{C1}): \emph{(a)} the accessor returns
$\mathrm{Option}$, so $\mathrm{None}$ (never held) and $\mathrm{Some}(0_P)$ (held, now
flat) stay distinct --- variation-margin settlement, wash-sale lookback, and record-date
entitlement each act on the difference; \emph{(b)} the carrier is monotone --- a row, once
created, is never removed, and close-out leaves it flat. The two are independent: (a)
constrains the return type, (b) the domain. The monotone carrier keeps the key set stable
across any checkpoint cut, which is what makes replay a literal fold (P8). The
high-water-mark and entry-NAV writers are valuation events specified in the valuation
satellite (\S\ref{sec:valuation}); they are out of scope here and leave those fields at
their defaults throughout.

\begin{lstlisting}[language=Haskell]
data PositionState = PositionState
  { psAc       :: !Qty           -- conserved, additive        <- Settle / Trade
  , psBalance  :: !Qty           -- conserved, additive        <- Transfer
  , psHwm      :: !Qty           -- NOT conserved, monotone    <- FeeCrystallise (qmax)
  , psEntryNav :: !(Maybe Qty)   -- NOT conserved, write-once  <- Subscribe
  } deriving (Eq, Show)

zeroP :: PositionState           -- flat row on first touch; also the held-and-flat baseline
zeroP = PositionState mempty mempty mempty Nothing

position :: Ledger -> WalletId -> UnitId -> Maybe PositionState   -- C1 Option accessor:
position l w u = Map.lookup (w, u) (ledgerPS l)
  -- Nothing = never held;  Just zeroP = held once, now flat.  The two stay distinct.
  -- Monotone carrier (C1, storage half): applyTx only inserts/updates a row, and no
  -- row deleter is exported anywhere -- close-out leaves a flat row, never removes it.
\end{lstlisting}

\paragraph{The three homes, three maps --- the sealed Ledger.}
Terms and status share the unit key but keep separate maps --- \texttt{ledgerPT}
append-only, \texttt{ledgerUS} overwrite-in-place --- because their mutation disciplines
differ. Their co-presence, ``$u$ registered in terms $\iff$ $u$ registered in status,'' is
preserved by construction, not policed: \texttt{register} is the sole introducer and writes
both, and every later writer touches them only for an already-registered unit. The position
map is the balance map enriched. The \texttt{Ledger} constructor and its field selectors are
unexported: a \texttt{Ledger} is built only by the writers below and read only through the
accessors. Were a selector exported, record update through it would install a
non-conserving map without naming the constructor, bypassing the single-writer discipline.
Coherence is carried by the single introducer; the seal keeps conservation true by
construction.

\begin{lstlisting}[language=Haskell]
data Ledger = Ledger
  { ledgerPT :: Map UnitId             ProductTerms   -- append-only terms, by unit
  , ledgerUS :: Map UnitId             UnitStatus     -- overwrite status,  by unit
  , ledgerPS :: Map (WalletId, UnitId) PositionState  -- positions,   by (holder, unit)
  }
  -- exported ABSTRACT: constructor and selectors unexported. A Ledger is built only by the
  -- writers below and read only through the accessors -> single-writer discipline; with no
  -- row deleter, the carrier is monotone (C1).

emptyLedger :: Ledger
emptyLedger = Ledger Map.empty Map.empty Map.empty

data LedgerError = ReRegistration UnitId | UnitAlreadyExists UnitId | UnknownUnit UnitId
                 | DanglingSupersede UnitId  -- SetSupersededBy target is not a registered unit
  deriving (Eq, Show)
\end{lstlisting}

\paragraph{The Transaction --- the one atomic event (C2, C3).}
One event records, for a single unit, the whole change across the three maps, in two
complementary parts. First, the conserved economic flow in \emph{edge} form: a list of
moves, each debiting one wallet and crediting another by the same magnitude, so the signed
per-unit sum of any move list is the group identity \emph{by construction} ---
$\mathrm{qneg}\,q \mathrel{\diamond} q = \mathit{mempty}$ for each leg, summed by the monoid
laws, the empty list (a move-less registration or amendment) included, which is the vacuous
base (\hyperref[cond:C9]{C9}). Second, the non-balance state delta that rides alongside the
moves: per-wallet bookkeeping (accumulated cost, high-water mark, entry NAV, via the C11
\texttt{FieldWrite} table), the shared \emph{UnitStatus} writes, and a \emph{ProductTerms}
introduction or append. Conservation is then a property of the type, not a checked
side-condition: there is no unconserved \texttt{Transaction} to reject, hence no
\texttt{validate} gate (\hyperref[cond:C2]{C2}). It applies in full or not at all; a
partially applied transaction is not a representable state (\hyperref[cond:C3]{C3}). One
structural field, \texttt{txIntroduce}, is the sole asymmetry between registration and every
other event: \texttt{Just} introduces the unit (writing terms and status together),
\texttt{Nothing} operates on an existing one --- so registration is the move-less instance of
this one type, not a separate operation. An event touching several units --- a rebalance
trading three futures, or a future bought against cash --- is a transaction per unit applied
as a single fold; per-unit conservation composes to event-level conservation. Induction over
the stream lifts the per-event identity to the global invariant P1.

\begin{lstlisting}[language=Haskell]
-- One value, two notations: txMoves is the per-edge view of the conserved flow; the
-- remaining fields are the non-balance bookkeeping that rides with it.
-- (Move, moveDelta, mUnit, mQty are the framework primitives of sec:ledger.)
data Transaction = Transaction
  { txUnit      :: UnitId                           -- the unit this event concerns
  , txMoves     :: [Move]                           -- conserved flow as edges (by construction)
  , txRows      :: Map WalletId [SomeWrite]         -- non-balance per-wallet bookkeeping
  , txStatus    :: [StatusWrite]                    -- shared UnitStatus writes
  , txIntroduce :: Maybe (TermsVersion, UnitStatus) -- Just => introduce txUnit (registration)
  , txAppend    :: Maybe TermsVersion               -- Just => append a version (C6, Preserving)
  } deriving (Show)

-- Conservation witnesses. Each move contributes qneg q <> q = mempty, so the foldMap lands
-- at mempty BY THE MONOID LAWS -- not by a sum computed and compared at run time (C2, C9).
unitDelta :: UnitId -> Transaction -> Qty
unitDelta u tx =
  foldMap (\m -> if mUnit m == u then qneg (mQty m) <> mQty m else mempty) (txMoves tx)
-- law:  unitDelta u tau == mempty   for every unit u and transaction tau

netDelta :: Transaction -> Qty                       -- the empty/registration case too
netDelta tx = foldMap (\m -> qneg (mQty m) <> mQty m) (txMoves tx)

-- The canonical move-less instances. Trade/settlement/transfer use the record constructor
-- directly (txIntroduce = Nothing, moves + rows present).
registerTx :: UnitId -> TermsVersion -> UnitStatus -> Transaction   -- C5+C7+C10: introduction
registerTx u tv us = Transaction
  { txUnit = u, txMoves = [], txRows = Map.empty
  , txStatus = [], txIntroduce = Just (tv, us), txAppend = Nothing }

appendTx :: UnitId -> TermsVersion -> Transaction                   -- C6 Preserving amendment
appendTx u tv = Transaction
  { txUnit = u, txMoves = [], txRows = Map.empty
  , txStatus = [], txIntroduce = Nothing, txAppend = Just tv }

supersedeTx :: UnitId -> UnitId -> Transaction                      -- C8 Breaking: stamp old unit
supersedeTx u uFresh = Transaction
  { txUnit = u, txMoves = [], txRows = Map.empty
  , txStatus = [SetSupersededBy uFresh], txIntroduce = Nothing, txAppend = Nothing }
\end{lstlisting}

\paragraph{Per-field canonical writer (C11).}
Each \emph{PositionState} field is tagged with the closed set of field-writers permitted to
mutate it: accumulated cost by settle and trade; high-water mark by fee-crystallisation;
entry NAV by subscription. The balance field is not in this table: its sole canonical writer
is the \texttt{Move} edge (handler \texttt{Transfer}), applied in \texttt{applyTx}, so a
balance change without a conserving move is unrepresentable. A write by a field-writer
outside the permitted set
is a type error at the writer's authorship site; the tag is erased once writes share one
delta row, so the guarantee binds at authorship, not at the stored row
(\hyperref[cond:C11]{C11}). These field-writers are a different axis from the event classes
of \hyperref[cond:C2]{C2}: one event class drives one or more field-writers, and the two
name-sets are not meant to coincide. The discipline is carried by a GADT (a generalised
algebraic data type, in which each constructor fixes its own type index) whose
constructors \emph{are} the field-to-writer table.

\begin{lstlisting}[language=Haskell]
data Handler = Settle | Trade | Transfer | FeeCrystallise | Subscribe deriving (Eq, Show)

-- The GADT constructors ARE the C11 field->writer table; no other pairing is representable.
-- The balance field is absent: its sole writer is the Move edge (handler 'Transfer), applied
-- in applyTx -- so "balance written off-ledger" is unrepresentable.
data FieldWrite (h :: Handler) where
  WAc       :: Qty -> FieldWrite 'Settle           -- ac, by settle
  WAcTrade  :: Qty -> FieldWrite 'Trade            -- ac, by trade
  WHwm      :: Qty -> FieldWrite 'FeeCrystallise   -- hwm, by fee crystallisation
  WEntryNav :: Qty -> FieldWrite 'Subscribe        -- entry_nav, by subscribe
deriving instance Show (FieldWrite h)

data SomeWrite where SomeWrite :: FieldWrite h -> SomeWrite  -- index erased once writes share a row
instance Show SomeWrite where show (SomeWrite w) = show w

applyWrite :: FieldWrite h -> PositionState -> PositionState  -- total over all constructors
applyWrite (WAc q)       p = p { psAc      = psAc p <> q }
applyWrite (WAcTrade q)  p = p { psAc      = psAc p <> q }
applyWrite (WHwm q)      p = p { psHwm      = qmax (psHwm p) q }                    -- monotone
applyWrite (WEntryNav q) p = p { psEntryNav = maybe (Just q) Just (psEntryNav p) } -- write-once

-- A settle handler authors ONLY 'Settle writes -- the result type IS the C11 checkpoint.
settleHandler :: [(WalletId, Qty)] -> Map WalletId [FieldWrite 'Settle]
settleHandler legs = Map.fromList [ (w, [WAc q]) | (w, q) <- legs ]

erase :: Map WalletId [FieldWrite h] -> Map WalletId [SomeWrite]  -- authorship -> erasure (S3)
erase = fmap (map SomeWrite)

_c11_ok_settle :: Qty -> FieldWrite 'Settle
_c11_ok_settle = WAc
_c11_ok_fee    :: Qty -> FieldWrite 'FeeCrystallise
_c11_ok_fee    = WHwm
-- _c11_bad :: Qty -> FieldWrite 'Settle
-- _c11_bad = WHwm    -- TYPE ERROR: WHwm :: FieldWrite 'FeeCrystallise, not 'Settle
\end{lstlisting}

\paragraph{Registration, settlement, moves, and amendment --- one door, \texttt{applyTx}.}
Every change to the three maps flows through one transition, \texttt{applyTx}; the ledger is
sealed and there is no second write path. Registration is the move-less transaction whose
\texttt{txIntroduce} is \texttt{Just}: it writes its first terms version and default status
as one entry, and refuses a unit already present --- re-registration is a hard error, never a
silent reset (\hyperref[cond:C10]{C10}). Because terms and status are written together,
``registered in terms $\iff$ registered in status'' holds by construction. \texttt{register}
is therefore \texttt{applyTx} of \texttt{registerTx}, not a separate operation. A move's
effect is its per-wallet net delta, written into the balance field and gated on registration;
its two legs are exact inverses, so it sums to zero over wallets, and a zero-quantity or
self-move nets to zero on every wallet and writes no net change --- \emph{held} means named in
a move that nets nonzero. Amendment has two tracks, chosen by a product-declared total
predicate (\hyperref[cond:C8]{C8}): a fungibility-preserving amendment is one move-less
transaction appending a terms version, leaving positions untouched; a fungibility-breaking
amendment spans two units, so it is the composition of two transactions --- introduce the
successor $u_{\mathrm{new}}$, then stamp
\texttt{superseded\_by}$(u_{\mathrm{old}}\!\to u_{\mathrm{new}})$ on the predecessor's status
--- run atomically through \texttt{applyTx} by \texttt{amend}, the move of holders being a
separate paired-issuance event, never part of the amendment. Ordering the two transactions so
the successor is introduced before the predecessor is stamped means the supersede pointer
always names a registered unit; a \texttt{SetSupersededBy} whose target is absent is rejected
as a \texttt{DanglingSupersede}, so the pointer can never dangle. Reads are capability-scoped: a
cross-overlay read of another holder's position is forbidden, and strategy exports flow only
through \emph{UnitStatus} (\hyperref[cond:C4]{C4}).

A move carries its conserved flow as edges in \texttt{txMoves}; its two legs are exact
inverses (\texttt{qneg}), so the signed per-unit sum is \texttt{mempty} by construction, and
a zero-quantity or self-move contributes no net write.

\begin{lstlisting}[language=Haskell]
-- C5+C7+C10: register is applyTx of the move-less introduction. It is not a primitive
-- standing outside the fold; registration, trade, settlement and amendment all flow through
-- the ONE apply.
register :: UnitId -> TermsVersion -> UnitStatus -> Ledger -> Either LedgerError Ledger
register u tv us = applyTx (registerTx u tv us)

-- The single door. C3 (all-or-nothing) is structural: the WHOLE new Ledger or a Left, never
-- a partial state. txIntroduce=Just must introduce an ABSENT unit (else ReRegistration, C10),
-- writing PT and US together (C5/C7); txUnit and every unit a move names must be registered
-- (else UnknownUnit), so PT<->US co-presence is preserved; rows insert/update, never delete
-- (C1 monotone). Conservation is already a property of txMoves (edges) -- nothing to validate.
applyTx :: Transaction -> Ledger -> Either LedgerError Ledger
applyTx tx l0 = do
    l1 <- introduce l0
    ()  <- requireKnown l1
    ()  <- requireSupersedeTargets l1
    Right (commit l1)
  where
    u = txUnit tx

    introduce l = case txIntroduce tx of
      Nothing -> Right l
      Just (tv, us)
        | u `Map.member` ledgerPT l -> Left (ReRegistration u)              -- C10
        | otherwise -> Right l
            { ledgerPT = Map.insert u (ProductTerms (tv :| [])) (ledgerPT l) -- C7
            , ledgerUS = Map.insert u us                        (ledgerUS l) -- C5
            }

    requireKnown l =
      case filter (\x -> not (x `Map.member` ledgerPT l)) (u : map mUnit (txMoves tx)) of
        (bad : _) -> Left (UnknownUnit bad)
        []        -> Right ()

    -- Referential integrity for the supersede pointer: a SetSupersededBy target must name a
    -- registered unit, else the stamp would dangle. A Breaking amendment introduces the
    -- successor first, so its id is registered by the time supersedeTx is applied (C8).
    requireSupersedeTargets l =
      case filter (\x -> not (x `Map.member` ledgerPT l))
                  [ t | SetSupersededBy t <- txStatus tx ] of
        (bad : _) -> Left (DanglingSupersede bad)
        []        -> Right ()

    commit l = l
      { ledgerPT = maybe (ledgerPT l)
                         (\tv -> Map.adjust (appendVersion tv) u (ledgerPT l))   -- C6 append
                         (txAppend tx)
      , ledgerUS = if null (txStatus tx)
                     then ledgerUS l
                     else Map.adjust (\s -> foldl' (flip applyStatus) s (txStatus tx)) u (ledgerUS l)
      , ledgerPS = foldl' (flip applyMoveRow)
                          (Map.foldrWithKey applyRow (ledgerPS l) (txRows tx))
                          (txMoves tx)
      }

    applyRow w writes ps =                                      -- non-balance bookkeeping
      let key = (w, u)
          cur = Map.findWithDefault zeroP key ps               -- first touch -> flat row
          new = foldl' (\acc (SomeWrite fw) -> applyWrite fw acc) cur writes
      in  Map.insert key new ps                                 -- insert/update only

-- A Move writes ONLY psBalance, via the conserving pair from moveDelta (from -= q, to += q),
-- keyed by the move's own unit -- the sole writer of the balance field.
applyMoveRow :: Move -> Map (WalletId, UnitId) PositionState
             -> Map (WalletId, UnitId) PositionState
applyMoveRow m ps0 = foldl' step ps0 (moveDelta m)
  where
    step ps (w, dq) =
      let key = (w, mUnit m)
          cur = Map.findWithDefault zeroP key ps
      in  Map.insert key cur { psBalance = psBalance cur <> dq } ps

-- C8: a product-declared total predicate chooses the track. Either track reaches the maps
-- ONLY through applyTx -- there is no second write path.
data Fungibility = Preserving | Breaking deriving (Eq, Show)
type FungibilityPredicate = ProductTerms -> TermsVersion -> Fungibility
data AmendResult = Appended UnitId | Superseded UnitId UnitId deriving (Eq, Show)

amend :: FungibilityPredicate -> UnitId -> TermsVersion -> UnitId -> Ledger
      -> Either LedgerError (AmendResult, Ledger)
amend isFungible u tvNew uFresh l =
  case Map.lookup u (ledgerPT l) of
    Nothing -> Left (UnknownUnit u)
    Just pt -> case isFungible pt tvNew of
      Preserving ->                                 -- ONE move-less Transaction; positions untouched
        (\l' -> (Appended u, l')) <$> applyTx (appendTx u tvNew) l
      Breaking                                      -- spans two units: a composite of TWO Transactions
        | uFresh `Map.member` ledgerPT l -> Left (UnitAlreadyExists uFresh)
        | otherwise -> do
            l1 <- applyTx (registerTx uFresh tvNew defaultStatus) l   -- introduce the successor
            l2 <- applyTx (supersedeTx u uFresh)                 l1   -- stamp the predecessor
            Right (Superseded u uFresh, l2)         -- old terms never rewritten (C7/P7)
\end{lstlisting}

\paragraph{Deterministic replay.}
Replay is the monadic left-fold of the event stream over the empty ledger. It is
deterministic because application is a pure, total function of the event and the prior
ledger: the same events give the same ledger. It returns the rebuilt ledger on a
well-formed stream and an error on an ill-formed one --- a repeated registration, or a move
or settlement on an unregistered unit --- the fold halting at the first refusal.
Checkpointing is sound by the monadic left-fold law: the fold over a concatenation splits
at any cut. Since registration, settlement, and moves are all events, replay from the empty
ledger rebuilds every unit's terms, status, and positions, so every view --- balances,
profit and loss, the cached \emph{UnitStatus} of \S\ref{sec:sec04-unitstatus} --- is a
projection of the stream. Row retention serves audit, not determinism (P8).

\begin{lstlisting}[language=Haskell]
-- A Transaction IS the event; applyTx is the pure, total transition. Replay is the
-- error-stopping left fold of the stream. Determinism (P8) is the fold-homomorphism law,
--     replay (xs <> ys)  =  replay xs >=> replay ys,
-- so the result is independent of where a checkpoint is cut. On a well-formed stream (every
-- event on a registered unit) no Left arises, so replay is total there. (Balances and PnL are
-- further projections of the same stream -- netBal, in the balance-sheet section.)
replay :: [Transaction] -> Ledger -> Either LedgerError Ledger
replay ds l0 = foldM (\l d -> applyTx d l) l0 ds
\end{lstlisting}

\subsection{The twelve conditions}
\label{sec:sec04-conditions}

Each condition renders one illegal state unrepresentable. The labels are stable tags, not
a sequence. Each is stated once below and forced at the construction step or instrument
named; the four instruments that force them --- a future whose state depends on past
trades, a managed account, a strategy that trades futures, and an instrument registered but
never traded --- are worked in \S\ref{sec:sec04-construction} and \S\ref{sec:sec04-three}.

\begin{longtable}{@{}p{0.8cm} p{11.0cm} p{2.5cm}@{}}
\toprule
 & \textbf{Condition} & \textbf{Forced at} \\
\midrule
\endhead
\hypertarget{cond:C1}{C1}\label{cond:C1} & Position accessor returns $\mathrm{Option}$
(never-held vs.\ held-and-flat distinct) \emph{and} the carrier is monotone (no row
deleted); both. & PositionState \\
\addlinespace
\hypertarget{cond:C2}{C2}\label{cond:C2} & Every \texttt{Transaction} carries its conserved
flow as edge moves, whose signed per-unit sum is \texttt{mempty} by construction; vacuous
base case (empty move list) included. & Transaction \\
\addlinespace
\hypertarget{cond:C3}{C3}\label{cond:C3} & A \texttt{Transaction} spans all three maps for
one unit and applies in full or not at all; partial application unrepresentable. &
Transaction \\
\addlinespace
\hypertarget{cond:C4}{C4}\label{cond:C4} & Reads are capability-scoped; cross-overlay reads
forbidden; strategy exports flow only through \emph{UnitStatus}. & strategy \\
\addlinespace
\hypertarget{cond:C5}{C5}\label{cond:C5} & \emph{UnitStatus} is registration-total:
product-declared defaults written at registration. & untraded unit \\
\addlinespace
\hypertarget{cond:C6}{C6}\label{cond:C6} & \emph{ProductTerms} is a non-empty version list,
append-only; no version mutated in place. & ProductTerms \\
\addlinespace
\hyperref[cond:C7]{C7} & \emph{ProductTerms} is registration-total: first
version written at registration. & ProductTerms \\
\addlinespace
\hypertarget{cond:C8}{C8}\label{cond:C8} & Two-track amendment by a total fungibility
predicate: \texttt{Preserving} appends a version; \texttt{Breaking} allocates a fresh $u$
and records \texttt{superseded\_by}. & amendment \\
\addlinespace
\hypertarget{cond:C9}{C9}\label{cond:C9} & A move-less event has an empty move list, which
folds to \texttt{mempty}, so it is conserved vacuously; conservation sums edge legs and never
divides by a holder count. & untraded unit \\
\addlinespace
\hyperref[cond:C10]{C10} & Re-registering an existing unit is a hard
error, never a silent reset. & registration \\
\addlinespace
\hypertarget{cond:C11}{C11}\label{cond:C11} & Each \emph{PositionState} field has a
canonical writer set; a writer outside it is a type error at authorship, erased at the
delta row. & PositionState \\
\addlinespace
\hypertarget{cond:C12}{C12}\label{cond:C12} & All per-(wallet, mandate/strategy) economic
state lives at \emph{PositionState}$[w,u_{\mathrm{MA}}]$; the $W$-sector is empty by
schema. & managed account \\
\bottomrule
\end{longtable}

\subsection{Why exactly three maps}
\label{sec:sec04-pareto-three}

Three independent constraints each force one map, and removing any one breaks its
constraint. \emph{(i)} Distinct per-holder state under the same contract forces a
$(w,u)$-keyed map --- \emph{PositionState}. \emph{(ii)} Shared observables read identically
by every holder force a $u$-keyed map distinct from per-holder state --- \emph{UnitStatus}.
\emph{(iii)} The two mutation disciplines --- append-only versioned terms versus an
overwrite-in-place status cache --- force a third, append-only map distinct from status ---
\emph{ProductTerms}. A fourth map, the $W$-sector, adds a sector whose economic content is
empty (\hyperref[cond:C12]{C12}). Three maps are the minimum basis of the problem, not a
compromise.

\subsection{The unique design}
\label{sec:sec04-pareto}

The adopted design --- three maps under C1--C12, design~B --- is the unique Pareto-optimum
on testability, correctness, and simplicity under any correctness gate. Each rejected
alternative fails on one forcing point.

\begin{itemize}[leftmargin=1.4em]
\item \textbf{A} (per-unit state, plus per-$(w,u)$ for futures only) has no sector for
shared observables distinct from immutable terms, so settlement prices and terms share one
overwrite cell --- the two mutation disciplines of \S\ref{sec:sec04-pareto-three}(iii)
collapse.
\item \textbf{C} (a sentinel unit carrying wallet state) has a sentinel with no issuer and
no conservation partner, so $\sum_{w}h(\cdot)=0$ fails by fiat.
\item \textbf{D} (an explicit $W$-sector) adds a sector whose economic content is empty:
every candidate datum is $(w,u_{\mathrm{MA}})$-keyed (\hyperref[cond:C12]{C12}).
\item \textbf{E} (state indexed only over the held set) equals B on correctness but has no
implementation in available tooling and is far less simple; shippability decides, not
elegance.
\item \textbf{F} (a universe wallet $w_\star$ folding \emph{UnitStatus} into
\emph{PositionState}$[w_\star,u]$) re-introduces the per-unit/per-$(w,u)$ split as a runtime
predicate on the wallet axis, and conservation must then carve out $w_\star$ by hand ---
\emph{UnitStatus} is the type-level expression of what F hides as a wallet-id convention.
\end{itemize}

\subsection{Invariants made unrepresentable}
\label{sec:sec04-unrep}

Under design~B with C1--C12, seven illegal states are placed beyond reach rather than
merely tested; each line names the invariant of the hub (\S\ref{sec:invariants}) it
discharges, or, where the guarantee is a contributing mechanism rather than the invariant
itself, the hub invariant it underpins. The P-numbers are the canonical hub numbering of
\S\ref{sec:invariants}. The mechanism differs by invariant and is named in each line. ``Unrepresentable'' is precise: an abstract
type whose only constructor is a checked smart constructor keeps the illegal value from
reaching the operation that would act on it; a non-empty list makes the versionless state
untypable; a per-field tag is a type error at authorship; conservation is carried by the edge
form of a move, whose two legs cancel by construction. It is not a claim that every guarantee
is a type-level fact --- read-scoping, in particular, is enforced at the capability layer
wrapping the accessors, not in the data reference.

\begin{longtable}{@{}p{2.0cm} p{6.7cm} p{5.8cm}@{}}
\toprule
\textbf{Invariant} & \textbf{Mechanism} & \textbf{Discharged by} \\
\midrule
\endhead
P1 (conservation of quantity) & Conserved flow is carried as \texttt{Move} edges, each
naming both legs (\texttt{qneg}\,$q \mathrel{\diamond} q = \mathit{mempty}$), so the signed
per-unit sum of any transaction is \texttt{mempty} by construction; there is no unconserved
\texttt{Transaction} to reject, hence no \texttt{validate} gate. & C2 (edge-conserved flow),
C3 (atomic transaction) \\
\addlinespace
P8 (snapshot/replay determinism) & Replay is the \texttt{foldM} of the stream, a fold
homomorphism: replaying a concatenated log equals replaying each part and composing; the
monotone carrier keeps the key set stable across any cut. & C1(b) (monotone carrier) \\
\addlinespace
P6 (lifecycle idempotency) & Lifecycle stage is $u$-keyed in \emph{UnitStatus},
written by replacement, so re-applying a stage write is the identity; non-conserved fields
carry the same algebra (high-water mark by $\max$, entry NAV write-once). & C5
(registration-total status) \\
\addlinespace
Terms immutability (underpins P4) & Append-only non-empty version list; no in-place
mutation, no exported setter. & C6 (append-only), C7 (registration-total) \\
\addlinespace
Identity immutability (underpins P4) & Re-registration rejected; a breaking amendment
allocates $u_{\mathrm{new}}$ and never rewrites $u_{\mathrm{old}}$. & C8 (two-track),
C10 (no re-registration) \\
\addlinespace
P7 (virtual/real isolation; read-scoping) & Cross-overlay reads forbidden at the capability
layer wrapping the accessors, not in the data reference. & C4 (capability-scoped reads) \\
\addlinespace
Handler-field canon (underpins P9) & Canonical writer set per field; a writer outside it is
a type error at authorship, erased once writes share a delta row. & C11 (per-field writer) \\
\bottomrule
\end{longtable}

No other candidate design places more than three of the core invariants beyond reach.

\subsection{Reference implementation}
\label{sec:sec04-reference}

The reference is \texttt{reference/Ledger.hs}, woven through \S\ref{sec:sec04-construction} and
collected here. It is total and deterministic, and makes the illegal states unrepresentable
rather than merely checked: \texttt{Ledger} is sealed with no row deleter (C1 monotone
half) and an \texttt{Option} accessor (C1 Option half); a \texttt{Transaction} carries its
conserved flow as \texttt{Move} edges, so the signed per-unit sum is \texttt{mempty} by
construction and no \texttt{validate} gate is needed (C2, the empty move list giving C9);
\texttt{ProductTerms} is abstract over a non-empty list grown only by registration and
\texttt{appendVersion} (C6, C7); the \texttt{FieldWrite} GADT is the C11 field-to-writer
table; \texttt{applyTx} is the single door, returning the whole ledger or an error with no
partial state (C3); registration is the move-less \texttt{txIntroduce}=\texttt{Just}
instance, rejecting re-registration (C10) and writing terms and status together; and
\texttt{amend} is the two-track C8. Every \emph{UnitStatus} field is written only by
\texttt{applyStatus}, folded inside \texttt{applyTx}, which is the projection-of-log
discipline of \S\ref{sec:sec04-unitstatus} held by construction.

The pieces woven through \S\ref{sec:sec04-construction} constitute \texttt{reference/Ledger.hs} in
full, building on the \texttt{Move} algebra of \S\ref{sec:ledger}
(\texttt{move}, \texttt{moveDelta}, \texttt{Timestamp}, \texttt{SourceId}); the consolidated,
compile-ordered source is collected as the section's reference. Run on that source, the
following evaluation makes the behaviour concrete --- the comments are the printed results.

\begin{lstlisting}[language=Haskell]
expect :: Show e => Either e a -> IO a
expect = either (error . show) pure

main :: IO ()
main = do
  let uES = UnitId "CME-ES"; wB = WalletId "BUYER"; wS = WalletId "SELLER"
  l0 <- expect $ register uES (TermsVersion "ES-v1" (Map.fromList [("mult","50")]))
                              defaultStatus emptyLedger
  print (position   l0 wB uES)        -- Nothing : never held (C1 Option half)
  print (unitStatus l0    uES)        -- Just (UnitStatus Registered Nothing Nothing) (C5)
  -- a conserving trade: +1000 to BUYER as a Move EDGE (balance), with accumulated-cost rows
  -- authored at 'Settle then erased (C11). (move/Timestamp/SourceId are the sec:ledger types.)
  Just mv <- pure (move wS wB uES (Qty 1000) (Timestamp 0) (SourceId "demo"))
  let tradeTx = Transaction uES [mv]
                  (erase (settleHandler [(wB, Qty 1000), (wS, Qty (-1000))]))
                  [] Nothing Nothing
  print (netDelta tradeTx)            -- 0 : the move legs cancel BY CONSTRUCTION (C2)
  l1 <- expect $ applyTx tradeTx l0
  print (psBalance <$> position l1 wB uES) -- Just 1000 : balance, written by the Move edge
  print (psAc      <$> position l1 wB uES) -- Just 1000 : accumulated cost, via the row
  -- There is no unconserved Transaction to reject: a one-legged flow is not a Move, so it is
  -- not representable in txMoves -- conservation is a property of the type, not a validate gate.
  -- close-out: a Move back keeps a flat row (monotone carrier): held-and-flat /= never-held
  Just mvC <- pure (move wB wS uES (Qty 1000) (Timestamp 1) (SourceId "demo"))
  l2 <- expect $ applyTx (Transaction uES [mvC] Map.empty [] Nothing Nothing) l1
  print (psBalance <$> position l2 wB uES) -- Just 0 : retained, not Nothing
  print (register uES (TermsVersion "dup" Map.empty) defaultStatus l2)
                                      -- Left (ReRegistration (UnitId "CME-ES")) (C10)
\end{lstlisting}

\begin{remark}[FORMALIS clearance]
The listings above are excerpted verbatim from \texttt{reference/\allowbreak Ledger.hs}, the
FORMALIS-cleared reference core; they carry that clearance and introduce no new code.
Quantities are exact \texttt{Integer} minor units, a correctness choice that makes
determinism and conservation arithmetic facts. Three structural points match the cleared
code: the per-unit homes are \emph{three coupled maps} (\texttt{ledgerPT}, \texttt{ledgerUS},
\texttt{ledgerPS}) with co-presence preserved by construction --- matching both the code and
\S\ref{sec:sec04-pareto-three}; the lifecycle is the closed sum
\texttt{Registered}/\texttt{Active}/\texttt{Expired}/\texttt{Closed} with the last settlement
price a separate \texttt{Maybe} field written only through a logged \texttt{StatusWrite},
rather than carried on an \texttt{Active} constructor; and one atomic \texttt{Transaction}
carries both the conserved move edges and the non-balance state delta, so conservation holds
by construction and there is no separate validate step.
\end{remark}

The testing commitments for this construction --- handler mutation score 85--90\%,
lifecycle-guard 70--80\%, core overall $\geq80\%$, and a state-machine model bounded by a
finite generator universe --- are carried by the testing agenda
(\S\ref{sec:conclusion}, Appendix~B), not restated here.
